\documentclass[11pt,a4paper]{article}
\usepackage[utf8]{inputenc}
\usepackage{hyperref}
\usepackage{booktabs}
\usepackage{graphicx}
\usepackage{amsmath}
\usepackage{geometry}
\geometry{margin=2.5cm}

\title{\textbf{Which Consciousness Modules Actually Do Anything?\\
An Ablation Study of a Deployed Voice Assistant}}

\author{
  Diyorbek\\
  Independent Researcher\\
  Tashkent, Uzbekistan\\
  \texttt{github.com/diyorbek2848/neurokey-consciousness}
}

\date{May 2026 (v2 — updated evaluation)}

\begin{document}
\maketitle

\begin{abstract}
We present NeuroKey, a deployed real-time voice assistant built around the functional
consciousness indicators of Butlin et al.\ (2023), and an ablation study measuring which
of its consciousness modules demonstrably affect its output.

Of Butlin et al.'s 14 indicators we implement \textbf{4 fully} (GWT-1, GWT-2, GWT-3, AE-2),
\textbf{4 partially} (HOT-2, PP-1, AST-1, AE-1) and \textbf{6 not at all} (RPT-1, RPT-2,
GWT-4, HOT-1, HOT-3, HOT-4). We additionally implement modules from Integrated Information
Theory and Self-Model Theory; we note that Butlin et al.\ explicitly exclude IIT as
incompatible with their working assumption of computational functionalism, so these fall
outside their framework and are reported separately.

Our contribution is measurement rather than coverage. Disabling each of 47 orchestrated
modules in turn, we find that \textbf{26 produce a measurable contribution to the assembled
context and 21 produce none} --- the null set including the modules implementing the paper's
own GWT, FEP and predictive-coding commitments. A sensitivity analysis over six diverse
inputs finds \textbf{20 of 29 context sections byte-identical regardless of input}: 73\,\% of
the assembled ``consciousness context'' carries no input-dependent information. The affective
section emits the same valence for ``I am exhausted'' and ``great news!''.

We also retract our previously reported score of 0.825/1.0. That evaluator returns 0.825 on a
28-module subset and 0.821 on the full 290-module system, a 46-module difference it cannot
detect, because 10 of its 14 checks assert module presence rather than behaviour.

We report these negative results because they are, to our knowledge, the first per-module
ablation published for a system of this kind, and because the alternative --- a score that
cannot fail --- is not evidence. \textbf{We invite independent evaluation.}
\end{abstract}

\section{Introduction}

Butlin et al.\ (2023) provided the first systematic framework for evaluating AI systems
against 14 indicators derived from neuroscientific theories of consciousness.
Their analysis of existing systems found partial satisfaction of several indicators
but no system had been \emph{designed} to satisfy the full framework.

NeuroKey addresses this directly. Each Butlin indicator is implemented as a distinct,
testable module. The system has been developed and used as a personal voice assistant since November 2024,
providing real-world validation beyond laboratory conditions.

\textbf{Key contributions:}
\begin{itemize}
  \item First deployed system designed explicitly around Butlin et al.\ (2023)
  \item Automated evaluator — any researcher can reproduce the score with \texttt{python consciousness\_evaluator.py}
  \item Persistent consciousness state that survives process restart
  \item Sub-650\,ms latency in production
  \item Open source: \url{https://github.com/diyorbek2848/neurokey-consciousness}
\end{itemize}

\section{Architecture}

The system comprises modules organized around four theoretical pillars.

\subsection{Global Workspace Theory (GWT)}
\texttt{global\_workspace.py} implements Baars-style competition:
modules submit signals with salience scores;
an EMA-weighted competition selects the spotlight;
the winner is broadcast to all subscribers.
Hysteresis prevents rapid oscillation between competing signals.

\subsection{Integrated Information Theory — Proxy}
True IIT $\Phi$ is NP-hard to compute exactly.
\texttt{phi\_proxy.py} implements a 3-node TPM
(valence, arousal, coherence as binary state variables)
following the IIT\,4.0 formalism.
The \texttt{pyphi} library is supported for exact calculation when available.

\subsection{Free Energy Principle (FEP)}
\texttt{active\_inference.py} implements Friston-style active inference cycles:
perception (minimize prediction error) and action (sample preferred states).
Each conversational turn runs a full inference cycle, producing a strategy
and surprise estimate that feed into the consciousness state.

\subsection{Episodic Memory}
Following Tulving (1972), episodes encode four dimensions:
\textit{time} (Unix timestamp),
\textit{place} (application context / screen state),
\textit{emotion} (valence $\times$ arousal at encoding),
\textit{content} (user utterance + assistant response).
Memory persists in SQLite across process restarts,
with full-text search over 1400+ stored conversations.

\section{Butlin et al.\ (2023) Evaluation}

\begin{table}[h]
\centering
\small
\begin{tabular}{llll}
\toprule
\textbf{Indicator} & \textbf{Description} & \textbf{Status} & \textbf{Module} \\
\midrule
GWT-1 & Parallel specialised systems      & Full    & 47 modules \\
GWT-2 & Limited-capacity workspace        & Full    & \texttt{global\_workspace.py} \\
GWT-3 & Global broadcast                  & Full    & \texttt{global\_workspace.py} \\
AE-2  & Embodiment: output--input model   & Full    & \texttt{embodied\_grounding.py} \\
\midrule
HOT-2 & Metacognitive monitoring          & Partial & \texttt{meta\_cognition.py} \\
PP-1  & Predictive coding in input        & Partial & \texttt{predictive\_layer.py} \\
AST-1 & Predictive model of own attention & Partial & \texttt{attention\_schema.py} \\
AE-1  & Agency under competing goals      & Partial & \texttt{goal\_system.py} \\
\midrule
RPT-1 & Algorithmic recurrence in input   & None    & --- \\
RPT-2 & Integrated perceptual reprs       & None    & --- \\
GWT-4 & State-dependent attention         & None    & --- \\
HOT-1 & Generative / noisy perception     & None    & --- \\
HOT-3 & Belief update from metacognition  & None    & --- \\
HOT-4 & Sparse, smooth coding             & None    & --- \\
\bottomrule
\end{tabular}
\caption{NeuroKey against the 14 indicators of Butlin et al.\ (2023):
4 full, 4 partial, 6 absent. The four ``partial'' entries are modules that load
but contribute nothing measurable to output (Table~\ref{tab:ablation}).
Note that IIT is not among Butlin et al.'s sources; they exclude it explicitly.}
\end{table}

\begin{table}[h]
\centering
\small
\begin{tabular}{lrl}
\toprule
\textbf{Module disabled} & \textbf{Characters lost} & \textbf{Section lost} \\
\midrule
\texttt{permanent\_memory}        & 8\,834 & conversation history \\
\texttt{vector\_memory}           & 2\,153 & semantic recall \\
\texttt{autobiographical\_memory} & 1\,991 & life narrative \\
\texttt{language\_of\_thought}    & 1\,296 & inner speech \\
\texttt{episodic\_memory}         & 1\,224 & long-term memory \\
\texttt{goal\_system}             & 1\,151 & goals \\
\midrule
\multicolumn{3}{l}{\textbf{26 of 47 modules contribute; 21 contribute nothing; 0 crash.}} \\
\multicolumn{3}{l}{Null set includes \texttt{global\_workspace}, \texttt{active\_inference},} \\
\multicolumn{3}{l}{\texttt{valence\_state}, \texttt{meta\_cognition}, \texttt{predictive\_layer}.} \\
\bottomrule
\end{tabular}
\caption{Ablation study. Each module disabled in turn; change in assembled context
measured across six diverse prompts, each condition in a separate process.}
\label{tab:ablation}
\end{table}

\textbf{Retraction of the previously reported score.} Earlier versions of this paper reported
a mean of 0.825/1.0 (13/14 passing) from \texttt{consciousness\_evaluator.py}. We withdraw
that figure. The evaluator returns 0.825 on a 28-module subset and 0.821 on the full
290-module system --- a 46-module difference it does not detect --- because 10 of its 14
checks are existence assertions of the form
\texttt{if hasattr(gw,\,"compete"): score = 0.7}. A check that cannot fail conveys no
information, and most checks additionally award a floor of 0.3--0.5 for merely not raising an
exception. We retain the script as a smoke test.

\textbf{Self-evaluation caveat.} The ablation and sensitivity studies remain self-administered,
but unlike the score they can fail and did: 21 modules returned null contributions and 20 of
29 context sections proved constant. Independent replication is invited; the harness is a
single file.

\section{Comparison with Related Systems}

\begin{table}[h]
\centering
\small
\begin{tabular}{lcccc}
\toprule
\textbf{Property} & \textbf{GPT-4} & \textbf{LIDA} & \textbf{OpenCog} & \textbf{NeuroKey} \\
\midrule
GWT implemented        & Partial & Yes & Partial & Yes \\
IIT metric             & No      & No  & No      & Proxy + pyphi \\
FEP / active inference & No      & No  & No      & Yes \\
Tulving episodic memory & No     & Yes & Yes     & Yes \\
Deployed real-time     & Yes     & No  & No      & Yes \\
$<$1\,s latency        & Yes     & No  & No      & Yes \\
Formal Butlin eval     & Post-hoc & No & No     & Built-in \\
Open source            & No      & Partial & Yes & Yes \\
\bottomrule
\end{tabular}
\caption{Comparison with related systems.}
\end{table}

The closest prior work is LIDA \cite{franklin2007}, which implements GWT
in a cognitive architecture, and OpenCog \cite{goertzel2014}.
Neither is deployed as a real-time voice assistant,
and neither was evaluated against Butlin et al.\ (2023).

\section{Limitations}

\begin{itemize}
  \item \textbf{Self-evaluation bias:} Independent evaluation is required for academic validity.
        The evaluator is written by the same author as the implementations.
  \item \textbf{GWT\_3 gap:} Attention ignition threshold not fully implemented.
        Spotlight selection is salience-weighted but lacks a hard ignition boundary.
  \item \textbf{IIT proxy:} The 3-node TPM approximation is not equivalent to true IIT $\Phi$.
  \item \textbf{No phenomenal claims:} All claims are functional/computational only.
        This system makes no assertion about subjective experience.
\end{itemize}

\section{Reproducibility}

\begin{verbatim}
git clone https://github.com/diyorbek2848/neurokey-consciousness
pip install -r requirements.txt
python consciousness_evaluator.py
\end{verbatim}

Output is printed to terminal and saved to \texttt{report.json}.
Each indicator includes score, pass/fail status, and evidence string explaining
the specific test performed.

\section{Conclusion}

A single developer can build and deploy a system organised around the Butlin et al.\ (2023)
indicators, and can keep it running daily for eighteen months. What that developer cannot do
by inspection alone is tell which of its parts are working. We built 47 consciousness modules
and could not, before measuring, have said which of them mattered; the answer turned out to be
26, with 21 contributing nothing detectable to the system's output, and with 73\,\% of the
context presented to the language model invariant across inputs as different as grief and good
news.

We think this is the more useful contribution. A consciousness score that rises when modules
are deleted --- as ours does for \texttt{phenomenal\_binding} and \texttt{goal\_system} --- is
not a measurement. An ablation that names its own dead modules is. The concrete next steps
follow directly from it: implement GWT-4 so that attention selects modules by state rather
than concatenating all of them, replace the placeholder prediction-error counter with genuine
predictive coding (PP-1), and close the loop from metacognitive monitoring to action selection
(HOT-3).

All code is released open source. We invite the research community to replicate these
measurements, and in particular to apply the same ablation to systems making comparable
claims.

\begin{thebibliography}{9}

\bibitem{butlin2023}
Butlin, P., Long, R., Elmoznino, E., Bengio, Y., et al.\ (2023).
Consciousness in artificial intelligence: Insights from the science of consciousness.
\textit{arXiv:2308.08708}.

\bibitem{baars1988}
Baars, B.\ J.\ (1988).
\textit{A Cognitive Theory of Consciousness}.
Cambridge University Press.

\bibitem{tononi2004}
Tononi, G.\ (2004).
An information integration theory of consciousness.
\textit{BMC Neuroscience, 5}(1), 42.

\bibitem{friston2010}
Friston, K.\ (2010).
The free-energy principle: A unified brain theory?
\textit{Nature Reviews Neuroscience, 11}(2), 127--138.

\bibitem{tulving1972}
Tulving, E.\ (1972).
Episodic and semantic memory.
In \textit{Organization of Memory} (pp.\ 381--403). Academic Press.

\bibitem{metzinger2003}
Metzinger, T.\ (2003).
\textit{Being No One}. MIT Press.

\bibitem{franklin2007}
Franklin, S., \& Patterson, F.\ G.\ (2007).
The LIDA architecture: Adding new modes of learning.
\textit{Integrated Design and Process Technology}.

\bibitem{goertzel2014}
Goertzel, B.\ (2014).
Artificial general intelligence: concept, state of the art, and future prospects.
\textit{Journal of Artificial General Intelligence, 5}(1), 1--48.

\end{thebibliography}

\end{document}
